%======================================================
%  FICHE D'ACCOMPLISSEMENT — main.tex
%  Étudiant : Youssef Ech-Chahlaoui
%  Organisme : Ministère de la Transition Numérique
%  Compiler  : pdflatex main.tex (deux passes)
%======================================================
\documentclass[a4paper,11pt]{article}

%─────────────────────────────────────────────────────
%  PACKAGES
%─────────────────────────────────────────────────────
\usepackage[utf8]{inputenc}
\usepackage[T1]{fontenc}
\usepackage[french]{babel}
\usepackage{lmodern}
\usepackage{microtype}
\usepackage{pifont}        % \ding{} pour symboles

% Mise en page
\usepackage[top=2.5cm, bottom=2.5cm, left=2.8cm, right=2.8cm]{geometry}
\usepackage{fancyhdr}
\usepackage{parskip}

% Couleurs
\usepackage[dvipsnames,table]{xcolor}
\definecolor{InstitBlue}{RGB}{0,51,102}
\definecolor{InstitGrey}{RGB}{100,100,110}
\definecolor{AccentBlue}{RGB}{30,120,190}
\definecolor{LightBg}{RGB}{240,244,250}
\definecolor{GreenOk}{RGB}{0,130,65}
\definecolor{OrangeWip}{RGB}{200,100,0}

% Listes
\usepackage{enumitem}
\setlist[itemize,1]{label=\textcolor{AccentBlue}{\textbullet}, leftmargin=1.5em}
\setlist[itemize,2]{label=\textcolor{InstitGrey}{$\triangleright$}, leftmargin=2em}
\setlist[enumerate,1]{leftmargin=1.5em, label=\textcolor{AccentBlue}{\arabic*.}}

% Tableaux
\usepackage{booktabs}
\usepackage{array}
\usepackage{tabularx}
\usepackage{multirow}

% Graphiques
\usepackage{graphicx}
\usepackage{float}
\usepackage{tikz}
\usetikzlibrary{positioning}

% Titres de sections
\usepackage{titlesec}
\titleformat{\section}
  {\color{InstitBlue}\large\bfseries\sffamily}
  {\thesection.}{0.6em}{}
  [\color{AccentBlue}\vspace{1pt}\rule{\linewidth}{1pt}]
\titlespacing*{\section}{0pt}{14pt}{6pt}
\titleformat{\subsection}
  {\color{InstitGrey}\normalsize\bfseries\sffamily}
  {\thesubsection.}{0.5em}{}

% Boîtes colorées (mdframed)
\usepackage{mdframed}
\mdfdefinestyle{SprintBox}{
  linecolor=AccentBlue, linewidth=1pt,
  backgroundcolor=LightBg,
  innerleftmargin=10pt, innerrightmargin=10pt,
  innertopmargin=8pt, innerbottommargin=8pt,
  roundcorner=4pt
}
\mdfdefinestyle{InfoBox}{
  linecolor=InstitBlue, linewidth=2pt,
  topline=false, rightline=false, bottomline=false,
  backgroundcolor=LightBg,
  innerleftmargin=10pt, innerrightmargin=8pt,
  innertopmargin=7pt, innerbottommargin=7pt
}

% Hyperliens
\usepackage[
  colorlinks=true,
  linkcolor=InstitBlue,
  urlcolor=AccentBlue,
  pdftitle={Fiche d'Accomplissement --- Youssef Ech-Chahlaoui},
  pdfauthor={Youssef Ech-Chahlaoui}
]{hyperref}

%─────────────────────────────────────────────────────
%  EN-TÊTES / PIEDS DE PAGE
%─────────────────────────────────────────────────────
\pagestyle{fancy}
\fancyhf{}
\renewcommand{\headrulewidth}{0.4pt}
\renewcommand{\footrulewidth}{0.4pt}
\fancyhead[L]{\small\color{InstitGrey}\sffamily Fiche d'Accomplissement de Stage PFA}
\fancyhead[R]{\small\color{InstitGrey}\sffamily Youssef Ech-Chahlaoui}
\fancyfoot[C]{\small\color{InstitGrey}\thepage}
\fancyfoot[R]{\small\color{InstitGrey}\sffamily Minist\`ere de la Transition Num\'erique}

%─────────────────────────────────────────────────────
%  COMMANDES UTILITAIRES
%─────────────────────────────────────────────────────
% Badges de statut
\newcommand{\done}{\textcolor{GreenOk}{\bfseries\sffamily\small \ding{51}\ Livr\'e}}
\newcommand{\wip}{\textcolor{OrangeWip}{\bfseries\sffamily\small \ding{109}\ En cours}}
\newcommand{\planned}{\textcolor{InstitGrey}{\sffamily\small \ding{113}\ Pr\'evu}}

% Tag technologique inline
\newcommand{\tech}[1]{%
  \tikz[baseline=(t.base)]
    \node[rounded corners=3pt, fill=AccentBlue!12, draw=AccentBlue!40,
          inner xsep=4pt, inner ysep=2pt,
          font=\scriptsize\sffamily\bfseries, text=InstitBlue](t){#1};}

% Coche verte inline
\newcommand{\ok}{\textcolor{GreenOk}{\bfseries\ding{51}}}
\newcommand{\approxok}{\textcolor{OrangeWip}{\bfseries$\approx$}}

%======================================================
%  DÉBUT DU DOCUMENT
%======================================================
\begin{document}
\thispagestyle{empty}

%─────────────────────────────────────────────────────
%  PAGE DE GARDE
%─────────────────────────────────────────────────────
\begin{tikzpicture}[remember picture, overlay]
  \fill[InstitBlue] (current page.north west)
    rectangle ([yshift=-3.8cm] current page.north east);
  \fill[AccentBlue] ([yshift=-3.8cm] current page.north west)
    rectangle ([yshift=-4.1cm] current page.north east);
  \fill[InstitBlue!8] (current page.south west)
    rectangle ([yshift=1.8cm] current page.south east);
\end{tikzpicture}

\vspace*{0.5cm}
\begin{center}
  {\color{white}\sffamily\Large\bfseries
    \'{E}cole Nationale Sup\'{e}rieure d'Informatique et d'Analyse des Syst\`{e}mes (ENSIAS)}\\[0.25cm]
  {\color{white!85!InstitBlue}\sffamily\normalsize
    G\'{e}nie Logiciel \& Syst\`{e}mes d'Information --- 3\textsuperscript{e} ann\'{e}e}
\end{center}

\vspace{3.8cm}

\begin{center}
  {\sffamily\Huge\bfseries\color{InstitBlue} Fiche d'Accomplissement}\\[0.35cm]
  {\sffamily\large\color{InstitGrey} de Stage Professionnel de Fin d'Ann\'{e}e (PFA)}\\[1cm]

  \noindent\colorbox{InstitBlue!8}{%
    \parbox{13cm}{\centering\vspace{6pt}
      {\sffamily\large\bfseries\color{InstitBlue}
        D\'{e}veloppement d'une Plateforme de D\'{e}mat\'{e}rialisation\\
        des Proc\'{e}dures Administratives en Ligne}\\[4pt]
      {\sffamily\small\color{InstitGrey}
        Direction des Syst\`{e}mes d'Information et de l'Innovation Num\'{e}rique}
      \vspace{4pt}}}
\end{center}

\vspace{1.5cm}

\begin{center}
\renewcommand{\arraystretch}{1.5}
\begin{tabular}{>{\bfseries\sffamily\color{InstitGrey}}r
                >{\color{InstitBlue}}l}
  \'{E}tudiant(e)            & Youssef Ech-Chahlaoui \\
  Fili\`{e}re                & G\'{e}nie Logiciel \& SI --- 3\textsuperscript{e} ann\'{e}e \\
  Organisme d'accueil       & Minist\`{e}re de la Transition Num\'{e}rique, Rabat \\
  D\'{e}partement             & Direction des Syst\`{e}mes d'Information et de l'Innovation \\
  P\'{e}riode                 & 01 avril 2025 au 30 juin 2025 (3 mois) \\
  Encadrant professionnel   & M. Khalid Alami (Chef de projet DSI) \\
  Encadrant acad\'{e}mique   & M. [À compléter --- Encadrant ENSIAS] \\
  Ann\'{e}e universitaire    & 2024--2025 \\
\end{tabular}
\end{center}

\vfill

\begin{center}
  {\small\color{InstitGrey}\sffamily
    Document confidentiel --- Usage acad\'{e}mique uniquement --- Juillet 2025}
\end{center}

\clearpage
\pagestyle{fancy}

%─────────────────────────────────────────────────────
%  SECTION 1 — CONTEXTE DU PROJET
%─────────────────────────────────────────────────────
\section{Contexte du projet}
\label{sec:contexte}

\begin{mdframed}[style=InfoBox]
Ce stage s'est d\'{e}roul\'{e} au sein du \textbf{Minist\`{e}re de la Transition
Num\'{e}rique} du Royaume du Maroc, pr\'{e}cis\'{e}ment au niveau de la
\textbf{Direction des Syst\`{e}mes d'Information et de l'Innovation Num\'{e}rique (DSIIN)},
\`{a} Rabat.
Le projet intitul\'{e}
\textbf{\guillemotleft~D\'{e}veloppement d'une Plateforme de D\'{e}mat\'{e}rialisation
des Proc\'{e}dures Administratives en Ligne~\guillemotright}
s'inscrit dans la strat\'{e}gie nationale \textbf{Maroc Digital 2030}, visant
\`{a} num\'{e}riser l'ensemble des services publics et \`{a} simplifier l'acc\`{e}s
des citoyens \`{a} l'administration.
\end{mdframed}

\subsection{Probl\'{e}matique}

Le Minist\`{e}re g\'{e}rait encore plusieurs proc\'{e}dures administratives ---
demandes de certificats, autorisations, d\'{e}clarations --- via des formulaires
papier ou des e-mails non structur\'{e}s. Cela entra\^{i}nait des d\'{e}lais de
traitement pouvant atteindre \textbf{15 jours ouvrables}, une absence de
tra\c{c}abilit\'{e} et une forte charge de travail pour les agents.
L'objectif \'{e}tait de concevoir et de d\'{e}ployer une plateforme web s\'{e}curis\'{e}e
permettant aux citoyens de soumettre leurs demandes en ligne et aux agents de les
traiter, valider et archiver num\'{e}riquement.

\subsection{Objectifs du stage}

\begin{enumerate}
  \item Analyser les proc\'{e}dures administratives existantes et formaliser le cahier
        des charges fonctionnel avec les parties prenantes.
  \item Concevoir l'architecture technique de la plateforme (API REST + portail citoyen + back-office agent).
  \item D\'{e}velopper les modules de d\'{e}p\^{o}t de dossier, de suivi de traitement et de notification.
  \item Int\'{e}grer le syst\`{e}me d'authentification national \textbf{CIN-Auth} (OAuth2 / OpenID Connect).
  \item D\'{e}ployer la plateforme sur l'infrastructure cloud souverain du minist\`{e}re.
  \item R\'{e}diger la documentation technique, le guide utilisateur et le rapport PFA.
\end{enumerate}

\subsection{P\'{e}rim\`{e}tre et contraintes}

\begin{itemize}
  \item \textbf{P\'{e}rim\`{e}tre :} 4 proc\'{e}dures cibles num\'{e}ris\'{e}es, portail
        citoyen, back-office agent, module de notification (SMS + e-mail), tableaux
        de bord de suivi.
  \item \textbf{Contraintes techniques :} Stack impos\'{e}e par la DSI --- Spring Boot
        (back-end), Angular (front-end), PostgreSQL, serveurs Linux de l'\'{e}tat.
  \item \textbf{Contraintes r\'{e}glementaires :} Conformit\'{e} RGPD-Maroc, chiffrement
        des donn\'{e}es personnelles, accessibilit\'{e} RGAA niveau AA.
  \item \textbf{Contraintes temporelles :} Proc\'{e}dure 1 en production avant fin mai~2025~;
        livraison compl\`{e}te fin juin~2025.
\end{itemize}

\clearpage

%─────────────────────────────────────────────────────
%  SECTION 2 — TÂCHES RÉALISÉES
%─────────────────────────────────────────────────────
\section{T\^{a}ches r\'{e}alis\'{e}es}
\label{sec:taches}

{\small\color{InstitGrey}\itshape
Projet men\'{e} selon un cycle it\'{e}ratif inspir\'{e} de la m\'{e}thodologie Agile,
avec des sprints bimensuels et des revues hebdomadaires avec M.~Alami (encadrant DSI).}

\vspace{0.8em}

%--- Sprint 1
\begin{mdframed}[style=SprintBox]
{\sffamily\bfseries\color{InstitBlue}
  Sprint 1 --- Semaines 1--2 : Prise en main \& Analyse}
\begin{itemize}
  \item Int\'{e}gration \`{a} l'\'{e}quipe DSIIN, acc\`{e}s aux environnements
        (VPN, GitLab interne, serveur de d\'{e}veloppement)
  \item \'{E}tude de l'existant : analyse de 4 proc\'{e}dures cibles via entretiens
        avec 6 agents et le responsable RH
  \item Installation et configuration de l'environnement local
        (IntelliJ IDEA, DBeaver, Postman, accès VPN ministère)
  \item R\'{e}daction du cahier des charges fonctionnel~:
        28~exigences fonctionnelles, 8~non-fonctionnelles, valid\'{e}es par le chef de projet
  \item Mise en place du GitLab interne~: repos, branches, MR templates
\end{itemize}
\end{mdframed}

\vspace{0.6em}

%--- Sprint 2
\begin{mdframed}[style=SprintBox]
{\sffamily\bfseries\color{InstitBlue}
  Sprint 2 --- Semaines 3--4 : Conception \& Architecture}
\begin{itemize}
  \item Mod\'{e}lisation UML~: diagramme de cas d'utilisation (3 acteurs : Citoyen,
        Agent, Administrateur), diagramme de classes (14 entit\'{e}s),
        4 diagrammes de s\'{e}quence (parcours critiques)
  \item Conception du sch\'{e}ma relationnel PostgreSQL (18~tables, migrations Liquibase)
  \item D\'{e}finition de l'API REST~: 41~endpoints document\'{e}s OpenAPI~3.0 (Swagger)
  \item Maquettes UI/UX dans Figma~: 16~\'{e}crans (portail citoyen + back-office agent)
        valid\'{e}s par la DSI et un panel de 3~agents-testeurs
  \item Architecture d\'{e}finie~: Spring Boot~3.2 + Angular~17 + Keycloak (SSO)
        + PostgreSQL~15 + Nginx
\end{itemize}
\end{mdframed}

\vspace{0.6em}

%--- Sprint 3
\begin{mdframed}[style=SprintBox]
{\sffamily\bfseries\color{InstitBlue}
  Sprint 3 --- Semaines 5--7 : D\'{e}veloppement --- Portail Citoyen}
\begin{itemize}
  \item Back-end~: impl\'{e}mentation des services
        \texttt{DemandeService}, \texttt{DossierService}, \texttt{NotificationService}~;
        28~endpoints REST op\'{e}rationnels
  \item Base de donn\'{e}es~: migrations Liquibase V1--V8, jeux de donn\'{e}es de test
  \item Front-end Angular~: modules d\'{e}p\^{o}t de dossier, suivi de demande,
        historique, espace personnel citoyen
  \item Int\'{e}gration Keycloak~: authentification SSO, gestion des rôles
        (\texttt{CITOYEN}, \texttt{AGENT}, \texttt{ADMIN})
  \item Tests unitaires~: 94~tests JUnit5 / Mockito, couverture~83\%
  \item Pipeline GitLab CI~: \'{e}tapes build, test, SonarQube, packaging
\end{itemize}
\end{mdframed}

\vspace{0.6em}

%--- Sprint 4
\begin{mdframed}[style=SprintBox]
{\sffamily\bfseries\color{InstitBlue}
  Sprint 4 --- Semaines 8--10 : Back-Office Agent \& Notifications}
\begin{itemize}
  \item D\'{e}veloppement du back-office agent~: file de traitement des demandes,
        validation / rejet / demande de compl\'{e}ment, annotation et archivage num\'{e}rique
  \item Syst\`{e}me de notifications multicanal~:
        e-mail (JavaMail + templates Thymeleaf) + SMS (API \textit{MTN/Moroccan carrier})
  \item Tableau de bord analytique admin~: statistiques en temps r\'{e}el
        (demandes par proc\'{e}dure, d\'{e}lais moyens de traitement, taux de r\'{e}solution)
  \item Mise en œuvre du chiffrement des donn\'{e}es sensibles (AES-256)
        et journalisation des acc\`{e}s (audit log)
  \item Tests d'int\'{e}gration~: 31~tests Testcontainers (PostgreSQL + Keycloak r\'{e}els)
\end{itemize}
\end{mdframed}

\vspace{0.6em}

%--- Sprint 5
\begin{mdframed}[style=SprintBox]
{\sffamily\bfseries\color{InstitBlue}
  Sprint 5 --- Semaines 11--12 : Tests, D\'{e}ploiement \& Passation}
\begin{itemize}
  \item Tests E2E~: 12~sc\'{e}narios Cypress (d\'{e}p\^{o}t de dossier, validation agent,
        notifications, acc\`{e}s refus\'{e})
  \item Recette utilisateur~: 3~sessions avec agents r\'{e}els, 18~anomalies corrig\'{e}es
  \item D\'{e}ploiement~: configuration Ansible pour serveurs Linux de l'\'{e}tat,
        certificats TLS Let's Encrypt, Nginx reverse-proxy
  \item \textbf{Proc\'{e}dure 1 mise en production} le 26~mai~2025
        (certification de conformit\'{e} num\'{e}rique)
  \item Documentation~: guide technique (55~pages) + guide agent (22~pages) + guide citoyen (10~pages)
  \item Formation~: 2~sessions de 2h pour 8~agents de la DSIIN
\end{itemize}
\end{mdframed}

\clearpage

%─────────────────────────────────────────────────────
%  SECTION 3 — COMPÉTENCES TECHNIQUES
%─────────────────────────────────────────────────────
\section{Comp\'{e}tences techniques mobilis\'{e}es}
\label{sec:competences}

\subsection{Langages de programmation}

\begin{center}
\renewcommand{\arraystretch}{1.3}
\begin{tabularx}{\textwidth}{>{\bfseries\sffamily}l X c}
  \toprule
  \rowcolor{InstitBlue!10}
  \textbf{Langage} & \textbf{Usage dans le projet} & \textbf{Niveau} \\
  \midrule
  Java 17       & Back-end complet (services, contr\^{o}leurs, entit\'{e}s JPA, s\'{e}curit\'{e}) & Avanc\'{e} \\
  TypeScript    & Front-end Angular~17, composants, services, guards             & Interm\'{e}diaire \\
  SQL           & Requ\^{e}tes PostgreSQL, vues, triggers, migrations Liquibase   & Avanc\'{e} \\
  HTML / CSS    & Templates Angular, mise en page responsive, accessibilit\'{e} RGAA & Interm\'{e}diaire \\
  YAML          & Pipelines GitLab CI, configs Ansible, Keycloak realm            & Interm\'{e}diaire \\
  Bash          & Scripts de d\'{e}ploiement Ansible, sauvegardes BDD automatiques & D\'{e}butant \\
  \bottomrule
\end{tabularx}
\end{center}

\subsection{Frameworks \& biblioth\`{e}ques}

\begin{itemize}
  \item \tech{Spring Boot 3.2}   Socle back-end~: REST, JPA/Hibernate, Security, Validation, Mail
  \item \tech{Angular 17}        Architecture SPA~: modules standalone, RxJS, Angular Material
  \item \tech{Keycloak 24}       SSO / IAM~: OpenID Connect, RBAC, int\'{e}gration Spring Security
  \item \tech{Liquibase}         Versionnement et migration du sch\'{e}ma PostgreSQL
  \item \tech{Thymeleaf}         Templates e-mail HTML dynamiques (notifications)
  \item \tech{JUnit 5 / Mockito} Tests unitaires back-end (94 tests, 83\% de couverture)
  \item \tech{Testcontainers}    Tests d'int\'{e}gration avec PostgreSQL et Keycloak r\'{e}els
  \item \tech{Cypress}           Tests end-to-end des 12 parcours utilisateurs critiques
\end{itemize}

\subsection{Outils et environnements}

\begin{itemize}
  \item \textbf{Gestion de version :} \tech{Git} \tech{GitLab} (instance interne minist\`{e}re) ---
        workflow GitFlow (branches \texttt{feature/}, \texttt{develop}, \texttt{main})
  \item \textbf{CI/CD :} \tech{GitLab CI/CD} --- pipeline auto~:
        build \textrightarrow{} tests \textrightarrow{} SonarQube \textrightarrow{} packaging Docker \textrightarrow{} d\'{e}ploiement Ansible
  \item \textbf{Qualit\'{e} de code :} \tech{SonarQube} --- 0 issue bloquante en livraison finale
  \item \textbf{Conteneurisation :} \tech{Docker} \tech{Docker Compose} ---
        5 services (backend, frontend, keycloak, db, nginx)
  \item \textbf{D\'{e}ploiement :} \tech{Ansible} --- playbooks idempotents pour les serveurs Linux de l'\'{e}tat
  \item \textbf{Base de donn\'{e}es :} \tech{PostgreSQL 15} --- 18~tables, AES-256 sur donn\'{e}es personnelles
  \item \textbf{IDE :} \tech{IntelliJ IDEA 2024.1} \tech{VS Code}
  \item \textbf{Design UI :} \tech{Figma} --- 16~\'{e}crans maquett\'{e}s et valid\'{e}s
  \item \textbf{Gestion de projet :} \tech{GitLab Issues} --- stories, milestones, suivi du sprint
  \item \textbf{Documentation :} \tech{Confluence} (wiki technique) + \tech{Swagger UI}
\end{itemize}

\clearpage

%─────────────────────────────────────────────────────
%  SECTION 4 — RÉSULTATS / LIVRABLES
%─────────────────────────────────────────────────────
\section{R\'{e}sultats et livrables concrets}
\label{sec:resultats}

\begin{center}
\renewcommand{\arraystretch}{1.4}
\begin{tabularx}{\textwidth}{c >{\sffamily\bfseries}l X c}
  \toprule
  \rowcolor{InstitBlue!10}
  \textbf{\#} & \textbf{Livrable} & \textbf{Description} & \textbf{Statut} \\
  \midrule
  1 & API REST back-end
    & 41~endpoints Spring Boot + Keycloak, s\'{e}curis\'{e}s, document\'{e}s OpenAPI
    & \done \\
  2 & Portail citoyen Angular
    & SPA responsive 16~pages~: d\'{e}p\^{o}t, suivi, historique, profil
    & \done \\
  3 & Back-office agent
    & File de traitement, validation/rejet, archivage num\'{e}rique
    & \done \\
  4 & Module notifications
    & E-mail (Thymeleaf) + SMS multicanal, configurable par proc\'{e}dure
    & \done \\
  5 & Infrastructure de d\'{e}ploiement
    & Playbooks Ansible, Docker Compose prod, TLS, Nginx, sauvegarde auto BDD
    & \done \\
  6 & Suite de tests
    & 94~unitaires (83\%) + 31~int\'{e}gration + 12~E2E Cypress
    & \done \\
  7 & Documentation compl\`{e}te
    & Guide tech.~55p. + guide agent~22p. + guide citoyen~10p. (Confluence)
    & \done \\
  8 & Rapport de stage PFA
    & Rapport acad\'{e}mique complet (en cours de r\'{e}daction finale)
    & \wip \\
  \bottomrule
\end{tabularx}
\end{center}

\vspace{0.8em}

\subsection{Indicateurs de performance (KPIs)}

\begin{itemize}
  \item \textbf{Couverture de tests :} 83\,\% (objectif~: 75\,\%) \ok
  \item \textbf{Temps moyen de r\'{e}ponse API (P95) :} 210\,ms (objectif~: $<$\,400\,ms) \ok
  \item \textbf{User stories livr\'{e}es :} 26 sur 28 (2~report\'{e}es en V2) \approxok
  \item \textbf{Proc\'{e}dure mise en production avant la date cible :} \ok
  \item \textbf{R\'{e}duction du d\'{e}lai de traitement :}
        de 15~jours (papier) \`{a} 2~jours ouvrables (en ligne), soit $-87\,\%$ \ok
  \item \textbf{Agents form\'{e}s :} 8 sur 8 pr\'{e}vus \ok
  \item \textbf{Anomalies bloquantes en recette :} 0 \ok
\end{itemize}

%─────────────────────────────────────────────────────
%  SECTION 5 — DIFFICULTÉS ET SOLUTIONS
%─────────────────────────────────────────────────────
\section{Difficult\'{e}s rencontr\'{e}es et solutions apport\'{e}es}
\label{sec:difficultes}

\begin{center}
\renewcommand{\arraystretch}{1.45}
\begin{tabularx}{\textwidth}{>{\sffamily\bfseries\color{InstitBlue}}p{3.2cm}
                              X
                              >{\itshape}X}
  \toprule
  \rowcolor{InstitBlue!10}
  \textbf{Difficult\'{e}} & \textbf{Description} & \textbf{Solution apport\'{e}e} \\
  \midrule
  Int\'{e}gration Keycloak
    & Configuration SSO complexe~: acc\`{e}s aux endpoints prot\'{e}g\'{e}s en tests
      d'int\'{e}gration impossible avec Keycloak mock standard.
    & Utilisation de \textbf{Testcontainers} avec une vraie instance Keycloak~;
      provisionnement automatis\'{e} du realm et des clients via API REST. \\[4pt]

  Contraintes r\'{e}seau
    & L'infrastructure minist\'{e}rielle interdisait l'acc\`{e}s direct \`{a} Maven Central
      depuis les serveurs, bloquant le pipeline CI.
    & Mise en place d'un \textbf{proxy Nexus} local pour les d\'{e}pendances Maven et npm,
      configur\'{e} avec le soutien de l'\'{e}quipe syst\`{e}me. \\[4pt]

  Accessibilit\'{e} RGAA
    & Les composants Angular Material pr\'{e}sentaient des lacunes d'accessibilit\'{e}
      (contraste, navigation clavier, lecteurs d'\'{e}cran).
    & Surcharge des composants Material pour respecter le RGAA~AA~;
      audit WAVE + tests manuels NVDA. \\[4pt]

  Mont\'{e}e en comp\'{e}tence Angular
    & D\'{e}couverte d'Angular~17 (standalone components, signals) sans
      exp\'{e}rience pr\'{e}alable du framework.
    & Auto-formation intense (documentation officielle, Angular University),
      revues de code hebdomadaires avec l'encadrant. \\
  \bottomrule
\end{tabularx}
\end{center}

\vspace{1em}

\begin{mdframed}[style=InfoBox]
\textbf{Bilan personnel.}
Ce stage au Minist\`{e}re de la Transition Num\'{e}rique a \'{e}t\'{e} une exp\'{e}rience
professionnelle exceptionnelle~: contribuer concr\`{e}tement \`{a} la strat\'{e}gie
\textit{Maroc Digital 2030}, voir une proc\'{e}dure administrative mise en production et
utilis\'{e}e par de vrais citoyens est une satisfaction sans \'{e}gale.
Au-del\`{a} des comp\'{e}tences techniques acquises (Angular, Keycloak, Ansible,
s\'{e}curisation des donn\'{e}es personnelles), j'ai d\'{e}couvert les sp\'{e}cificit\'{e}s
du d\'{e}veloppement dans un contexte institutionnel~: processus de validation plus formels,
contraintes r\'{e}glementaires fortes, n\'{e}cessit\'{e} d'une documentation exhaustive.
\end{mdframed}

%─────────────────────────────────────────────────────
%  SECTION 6 — CE QUI RESTE À FAIRE
%─────────────────────────────────────────────────────
\section{Ce qui reste \`{a} faire}
\label{sec:reste}

Les points suivants n'ont pas \'{e}t\'{e} compl\'{e}t\'{e}s avant la fin du stage
et sont transmis \`{a} l'\'{e}quipe DSI ou planifi\'{e}s en \textbf{Version 2}~:

\begin{itemize}
  \item \textbf{Proc\'{e}dures 2, 3 et 4 en production :}
        La proc\'{e}dure~1 est en ligne. Les trois suivantes sont d\'{e}velopp\'{e}es et
        en recette, mais leur mise en production est pr\'{e}vue en septembre~2025
        par l'\'{e}quipe DSI. \textit{Estimation restante~: 2~semaines de recette chacune.}

  \item \textbf{G\'{e}n\'{e}ration de r\'{e}c\'{e}piss\'{e}s PDF officiels :}
        Impression des accusés de réception signés \'{e}lectroniquement (JasperReports + PKCS\#12).
        Maquette valid\'{e}e, d\'{e}veloppement non d\'{e}marr\'{e}.
        \textit{Estim\'{e}~: 5~jours.}

  \item \textbf{Signature \'{e}lectronique des d\'{e}cisions :}
        Int\'{e}gration de la signature num\'{e}rique certifi\'{e}e (ANCFCC)
        pour les d\'{e}cisions administratives. Fonctionnalit\'{e} bloquante
        sur validation juridique toujours en cours.

  \item \textbf{Tests de charge syst\'{e}matiques :}
        Sc\'{e}narios JMeter pour 500 utilisateurs simultan\'{e}s.
        Premiers r\'{e}sultats corrects, non int\'{e}gr\'{e}s au pipeline CI
        par manque de temps.

  \item \textbf{Rapport PFA (en cours) :}
        R\'{e}daction du rapport acad\'{e}mique complet.
        Structure d\'{e}finie, chapitres~1 et~2 r\'{e}dig\'{e}s,
        chapitres~3--5 en cours.
        \textit{Date de remise pr\'{e}vue~: [à compléter].}
\end{itemize}

%─────────────────────────────────────────────────────
%  FIN DU DOCUMENT
%─────────────────────────────────────────────────────
\vfill
\begin{center}
  \rule{9cm}{0.4pt}\\[4pt]
  {\small\color{InstitGrey}\sffamily
    Fiche d'accomplissement --- Youssef Ech-Chahlaoui ---
    Stage PFA 2024--2025 --- Minist\`{e}re de la Transition Num\'{e}rique}
\end{center}

\end{document}
